\chapter{Electric Circuits}

\section{Ohm's Law}

\textbf{Electric current}~\(I\) is the flow of electrons through a conductor.
It is a scalar measured in amperes (\unit{\ampere}),
also informally called "amps".
Although the unit of current is amperes,
a device that measures current is an \textbf{ammeter}.

Current is the derivative of charge, i.e.,
\begin{equation} \label{eq:current-charge}
    I = \odv{q}{t}.
\end{equation}

Despite electrons being the carriers of current,
the direction of current is conventionally reverse of the direction
of electron flow,
as if positive charges are flowing;
this is called "\textbf{conventional current}".

Electric currents occur in response to a voltage,
and the voltage that drives a circuit is the \textbf{electromotive
force}~(\(\epsilon\)) or \textbf{emf}.

The relation between the voltage~\(V\) and current~\(I\) is given by
\textbf{Ohm's law}:
\begin{equation} \label{eq:ohms-law}
    V = I R,
\end{equation}
where \(R\) is the \textbf{resistance}, which is measured in ohms~(\unit{\ohm}).

\textbf{Resistors} are devices that produce resistance.
From a practical viewpoint,
electrical devices such as lights and computers may also be
considered resistors.

\section{Kirchoff's Laws}

\textbf{Kirchoff's laws}:
\begin{enumerate}
    \item \textbf{Kirchoff's junction rule}
        states that the total input current and total output current
        of a junction must be equal to each other:
        \begin{equation}
            I_\mathrm{in} = I_\mathrm{out}.
        \end{equation}
        This follows from the conservation of charge
        (\cpageref{sec:electric-charges}).
    \item \textbf{Kirchoff's loop rule} states that all the voltages
        in a closed circuit must add up to \(\qty{0}{\volt}\):
        \begin{equation} \label{eq:kirchoffs-loop-rule}
            \sum_i V_i = 0.
        \end{equation}
        This results from the fact that electric fields are conservative
        [\pgRef{sec:electric-potential}].
\end{enumerate}

\section{Power}

\textbf{Power}~(\(P\)) is the work done per unit of time
and has units of watts (\unit{\watt}),
which are equivalent to \unit{\joule \per \second}.

In circuits,
power is most commonly discussed in the context of heat dissipated
from resistors.
Volts are equal to \unit{\joule \per \coulomb},
and amperes are equal to \unit{\coulomb \per \second}.
Therefore, we can use dimensional analysis to find that
\begin{equation}
    \begin{aligned}
        I V
        &= \ab(\unit{\joule \per \coulomb})
        \ab(\unit{\coulomb\per \second})
        = \unit{\joule \per \second} \\
        &= P \\
        P &= I V.
    \end{aligned}
\end{equation}
We can apply Ohm's law~[\pgRef{eq:ohms-law}] to replace either
current or voltage.

\section{RC Circuits}

\textbf{RC circuits} are composed of a resistor and capacitor
hooked up to a voltage source.
They charge and and discharge gradually.
Practically, any time you charge or discharge a capacitor,
it is going to take time since wires have an inherent resistance.
(You could try superconductors, but I doubt you have access to them.)

From Kirchoff's loop rule~[\pgRef{eq:kirchoffs-loop-rule}],
Ohm's law~[\pgRef{eq:ohms-law}], and \pgRef{eq:capacitance-def} we know
\begin{equation}
    \epsilon - I R - \frac{q}{C} = 0.
\end{equation}
From \pgRef{eq:current-charge}, we get
\begin{equation}
    \begin{aligned}
        \epsilon - \odv{q}{t} R - \frac{q}{C} &= 0 \\
        \odv{q}{t} R + \frac{q}{C} &= \epsilon \\
        \odv{q}{t} + \frac{q}{R C} &= \epsilon \\
        \odv{q}{t} e^{\frac{t}{R C}}
        + \frac{q}{R C} e^{\frac{t}{R C}}
        &= \epsilon e^{\frac{t}{R C}} \\
        \odv{}{t} \ab[q e^{\frac{t}{R C}}] &= \epsilon e^{\frac{t}{R C}} \\
        q e^{\frac{t}{R C}} &= \int \epsilon e^{\frac{t}{R C}} \odif{t}
    \end{aligned}
\end{equation}
Awkwardly, we need an integration constant but \(C\) is already taken.
Let us use \(C_\mathrm{i}\) instead, giving us
\begin{equation}
    \begin{aligned}
        q e^{\frac{t}{R C}} &=  \epsilon R C e^{\frac{t}{R C}} + C_\mathrm{i} \\
        q &=  \epsilon R C + C_\mathrm{i} e^{- \frac{t}{R C}} \\
    \end{aligned}
\end{equation}
